% Figure: a boom hinged to a wall at its foot and held by an inclined
% cable to a higher wall point, carrying a load at its tip. Three forces
% act on the boom: the cable tension, the wall-pin reaction, and the load.
% The figure shows the assembly and the boom FBD with the pin reaction
% direction left as an unknown to be found by concurrence.
\begin{figure}[htbp]
\centering
\begin{tikzpicture}[
  member/.style={draw=engblack, line width=1.6pt},
  cable/.style={draw=engblue, line width=1.3pt},
  load/.style={draw=engred, -{Stealth}, very thick},
  react/.style={draw=engamber, -{Stealth}, very thick},
  scale=1.0,
]
% ===== Left: the assembly =====
\begin{scope}[xshift=0cm]
  % wall
  \draw[draw=engblack, very thick] (0,-0.2) -- (0,3.6);
  \foreach \y in {0.1,0.5,...,3.4} {\draw[draw=enggray, thin] (0,\y) -- (-0.28,\y-0.2);}
  \coordinate (O) at (0,0.6);        % pin at foot
  \coordinate (Tip) at (3.4,1.4);    % boom tip, boom inclined up
  \coordinate (Cw) at (0,3.0);       % cable anchor high on wall
  % boom
  \draw[member] (O) -- (Tip);
  % cable from high wall point to tip
  \draw[cable] (Cw) -- (Tip);
  % pin marker at O
  \fill[engblack] (O) circle (0.07); \draw[draw=engblack, fill=white] (O) circle (0.12);
  % cable anchor marker
  \fill[engblack] (Cw) circle (0.07); \draw[draw=engblack, fill=white] (Cw) circle (0.12);
  % tip pin
  \fill[engblack] (Tip) circle (0.06);
  \node[font=\scriptsize, anchor=east] at (O) {$A$};
  \node[font=\scriptsize, anchor=east] at (Cw) {$C$};
  \node[font=\scriptsize, anchor=south west] at (Tip) {$B$};
  % load at tip
  \draw[load] (Tip) -- ++(0,-1.4) node[anchor=north, font=\small, engred] {$W$};
  \node[font=\small, anchor=north] at (1.8,-0.55) {(a) the assembly};
\end{scope}
% ===== Right: FBD of the boom with the three forces =====
\begin{scope}[xshift=6.4cm]
  \coordinate (O2) at (0,0.6);
  \coordinate (Tip2) at (3.4,1.4);
  \draw[member] (O2) -- (Tip2);
  \fill[engblack] (O2) circle (0.06);
  \fill[engblack] (Tip2) circle (0.06);
  \node[font=\scriptsize, anchor=north east] at (O2) {$A$};
  \node[font=\scriptsize, anchor=south west] at (Tip2) {$B$};
  % cable tension at tip: up and toward the wall along BC
  \draw[cable, -{Stealth}, very thick] (Tip2) -- ++(-1.7,0.8)
    node[anchor=south, font=\scriptsize, engblue] {$T$};
  % load down at tip
  \draw[load] (Tip2) -- ++(0,-1.4) node[anchor=north, font=\small, engred] {$W$};
  % pin reaction at A, direction unknown (drawn provisionally)
  \draw[react] (O2) -- ++(1.3,0.7)
    node[anchor=south west, font=\scriptsize, engamber] {$R_A$};
  \node[font=\scriptsize, anchor=north, engamber] at (0.7,0.55) {(direction by concurrence)};
  \node[font=\small, anchor=north] at (1.8,-0.55) {(b) FBD of the boom};
\end{scope}
\end{tikzpicture}
\caption[A boom held by a cable with a wall pin]{A boom hinged to a wall
at its foot $A$, held up by a cable from a higher wall point $C$ to the
boom tip $B$, with a load $W$ hanging at $B$. (a) The assembly. (b) The
free-body diagram of the boom. Exactly three forces act on it: the cable
tension $T$ along $BC$, the hanging load $W$ down at $B$, and the pin
reaction $R_A$ at the foot. Because the cable tension and the load both
pass through $B$, their lines of action cross there, so concurrence forces
the pin reaction to pass through $B$ as well: $R_A$ acts along $AB$,
fixing the one direction the pin would otherwise leave open. The three
magnitudes then close a single force triangle.}
\label{fig:vol03ch01:boom-wall-pin}
\end{figure}
